% MLA-format report for COMP 363, Week 04: Classic vs. Karatsuba Multiplication
\documentclass[12pt,letterpaper]{article}

\usepackage[margin=1in, headheight=15pt, headsep=0.35in]{geometry}
\usepackage{newtxtext, newtxmath}   % Times New Roman, text and math
\usepackage{setspace}
\usepackage{fancyhdr}

% MLA: double spacing, 0.5in first-line indent, no extra space between paragraphs
\doublespacing
\setlength{\parindent}{0.5in}
\setlength{\parskip}{0pt}

% MLA running head: last name and page number, top right, on every page
\pagestyle{fancy}
\fancyhf{}
\rhead{Horan \thepage}
\renewcommand{\headrulewidth}{0pt}

\newcommand{\code}[1]{\texttt{#1}}

\begin{document}

\noindent Jackson Horan\\
Professor Irakliotis\\
COMP 363: Algorithms\\
24 September 2026

{\centering Classic vs.\ Karatsuba Multiplication: A Timing Study\par}

Both algorithms fit $T(n) = rT(n/c) + n^d$ with $c = 2$, since \code{split\_in\_half} halves each
factor. Classic makes four recursive calls ($r = 4$); Karatsuba makes three ($r = 3$). All other work is
linear ($d = 1$): \code{split\_in\_half}, \code{shift}, \code{add\_digits}, and \code{subtract\_digits}
each pass once over at most $n$ digits, a fixed number of times per call. Since $r > c^d = 2$, the Master
Theorem predicts $O(n^2)$ for classic and $O(n^{\log_2 3}) \approx O(n^{1.585})$ for Karatsuba
(Irakliotis, ``Recurrence'').

Counting only digit multiplications, $M(n) = 4M(n/2)$ and $K(n) = 3K(n/2)$ with $M(1) = K(1) = 1$, so
$M(n) = 4^{\log_2 n} = n^2$ and $K(n) = 3^{\log_2 n} = n^{\log_2 3}$. My counts matched exactly:
$262{,}144 = 512^2$ and $19{,}683 = 3^9$ at $n = 512$.

Classic was faster through $n = 16$; Karatsuba won from $n = 32$ on, taking 77.5~ms vs.\ 281.6~ms at
$n = 512$. On small inputs its extra additions and list copies outweigh the one multiplication it saves.
Doubling $n$ multiplied classic's time by 4.0--4.1 and Karatsuba's by 3.0--3.1, as $\Theta(n^2)$ and
$\Theta(n^{\log_2 3})$ predict.

At $n = 512$, Karatsuba did $13.3\times$ fewer multiplications but ran only $3.6\times$ faster, because
most of the time goes to linear list work, and Karatsuba does more of it per call. That gap held at about
3.6 from $n = 32$ on: a constant factor, not a different growth rate.

For $n \geq 8$, medians stayed within 12\% of the minimum; tiny inputs varied up to $2.3\times$. Noise
only adds time, so the minimum of seven trials is the most repeatable estimate.

\newpage

{\centering Works Cited\par}

% MLA: 0.5in hanging indent on each entry
\setlength{\parindent}{0pt}
\everypar{\hangindent=0.5in \hangafter=1}

Irakliotis, Leo. ``Assignment: Classic vs.\ Karatsuba Multiplication---A Timing Study.''
\emph{COMP 363: Algorithms}, Fall 2026.

{-}{-}{-}. ``Recurrence Analysis.'' \emph{COMP 363: Algorithms}, Fall 2026,
\code{master\_theorem.ipynb}.

\end{document}
